\documentstyle[%


righttag%


,11pt%


,amssymb%


,russian%               remove in Eng.


,izvru%                 Set izven% in Eng.


]{amsart}





%





\newcommand{\la}{\langle}


\newcommand{\ra}{\rangle}


\newcommand{\const}{\operatorname{const}}


\newcommand{\tts}[1]{}





\theoremstyle{plain}


\theorembodyfont{\it}


\newtheorem{lemnonum}{\hskip\parindent Лемма}


\renewcommand{\thelemnonum}{}





\theoremstyle{definition}


\theorembodyfont{\rm}


\newtheorem{cornonum}{\hskip\parindent Следствие}


\renewcommand{\thecornonum}{}





%\hoffset=-15mm%                  For Eng.


\setcounter{page}{27}


\god{2005}


\nomer{8~(519)} %                        Remove in Eng.


%\nomer{Vol.\,49, No.\,8}%               Set in Eng.


%\str{\pageref{begin}--\pageref{end}}%  Set in Eng.


\udk{514.152}





\author{\it В.А.\,Кыров}


% NB:   ^^ remove in Eng.


\title{Двуметрические пространства}





\date{}


%\pagestyle{myheadings}%         Set in Eng.


\pagestyle{plain} %              Remove in Eng.


\begin{document}


%\thispagestyle{plain}%          Set in Eng.


\maketitle


\label{begin}








В данной работе проводится построение двуметрических


двумерных многообразий. Строятся двуметрические функции, определяются


согласованные связности и группы изометрий. Приводится интерпретация


этих пространств в термодинамике.


\medskip








{\bf 1. Двуметрические плоскости.}


В самом широком смысле двуметрическая геометрия --- это геометрия с


двумя расстояниями. В этой геометрии паре точек ставится в соответствие не


одно число, как обычно, а два. В [1] проводится полная


классификация двуметрических двумерных феноменологически


симметричных геометрий, т.\,е.


геометрий со следующим свойством. Существуют многообразие $M$, являющееся


плотным подмногообразием


прямого произведения $N\times N$ некоторого двумерного многообразия $N$ на


себя, и гладкая невырожденная функция $f: M \to R^2$. Функцию $f$ будем


называть {\it метрической функцией или двуметрикой, или двуметрической


функцией}, ее компоненты будем обозначать $(f^1,f^2)$.


Существует такая гладкая функция шести переменных


$\varPhi:R^3 \to R^2$ с компонентами $(\varPhi_1,\varPhi_2)$,


причем $\mathrm{rang}\,\varPhi=2$,


что для любого кортежа из трех произвольных точек $\la xyz\ra$, каждая


пара из которого принадлежит множеству $M$, имеет место функциональное


уравнение


\begin{equation*}


\varPhi(f(xy),f(xz),f(yz))=0.


\end{equation*}





Существуют всего две двумерные феноменологически симметричные


двуметрические структуры. Это геометрии с


метрическими функциями, которые имеют следующие координатные представления:


\begin{alignat}{2}


f^1(xy)&=x^1 - y^1, &\quad f^2(xy)&=x^2 - y^2;\\


f^1(xy)&=(x^1 - y^1)x^2, &\quad f^2(xy)&=(x^1 - y^1)y^2,


\end{alignat}


где $(x^1,x^2)$ --- координаты точки $x$, а $(y^1,y^2)$ --- координаты точки $y$.


Пространство, в котором определена данная двуметрическая структура,


будет обозначаться $F^2$. Заметим, что области определения двуметрических


функций (1) и (2) совпадают со всем прямым произведением $F^2\times F^2$.


Обратим внимание на то, что множество $F^2$ совпадает


с аффинной плоскостью.





Рассмотрим две бесконечно близкие точки $y=(x^1,x^2)$ и $x=(x^1+dx^1,x^2+


dx^2)$, на которых метрические функции (1) и (2) принимают значения


\begin{alignat}{2}


f^1&=dx^1, &\quad f^2&=dx^2;\tag{1${}'$} \\


f^1&=(dx^1)(dx^2+x^2), &\quad f^2&=(dx^1)x^2.\tag{2${}'$}


\end{alignat}


Переходя в метрических функциях $(1')$, $(2')$ от специальных координат $(x^1,


x^2)$ к произвольным $(y^1,y^2)$ по формулам $x^1=\psi^1(y^1,y^2)$,


$x^2=\psi^2(y^1,y^2)$, получим


\begin{alignat*}{2}


f^1&=a_i^1dy^i, &\quad f^2&=a_i^2dy^i;\\


%


f^1&=(a_i^1dy^i)(a_i^2dy^i+\xi^2), &\quad


f^2&=a_i^1dy^i\xi^2,


\end{alignat*}


где $a^i_j=\partial \psi^i/\partial y^j$, $\xi^2=\psi^2$, $i,j=1,2$.





Преобразование $g$ двуметрического двумерного


пространства $F^2$ назовем {\it движением},


если оно сохраняет метрическую функцию $f$, т.е. двуметрику $(f^1,f^2)$.


В [1] показано, что по метрической функции


$f$ находится двухпараметрическая группа движений $G_{F^2}$, 


а по этой группе движений


восстанавливается метрическая функция $f$ с точностью до гладкой


функции $\varphi: R^2 \to R^2$.


Решая эту задачу для выше приведенных пространств, приходим к группам движений


$G_{F^2}$,


которые задаются следующими уравнениями:





для метрической функции (1)


$$


x^{1 \prime}=x^1+b,\quad x^{2 \prime}=x^2+a;


$$





для метрической функции (2)


\begin{equation}


x^{1 \prime}=e^{\alpha} x^1+b,\quad x^{2 \prime}=e^{-\alpha} x^2,


\end{equation}


где $\alpha$, $a$, $b$ --- произвольные постоянные.


Объединяя эти системы в одну, имеем


\begin{equation}


x^{1 \prime}=e^{\varepsilon\alpha} x^1+b,\quad


x^{2 \prime}=e^{-\varepsilon\alpha} x^2+(1 - \varepsilon)a,


\end{equation}


где $\varepsilon=0, 1$. Обратим вримание на то, что система (4)


зависит от двух независимых непрерывных параметров: $a$ и $b$ при 


$\varepsilon=0$; $\alpha$ и $\beta$ при $\varepsilon=1$. 


Двухточечными инвариантами группы (3) являются также


следующие функции:


\begin{gather}


\eta=\frac{f^1(xy) - f^2(xy)}{f^1(xy)}=\frac{x^2 - y^2}{x^2},\\


A=f^1(xy) - f^2(xy)=(x^1 - y^1)(x^2 - y^2).


\end{gather}





Двуметрическая геометрия с метрической функцией (2) допускает содержательную


интерпретацию в термодинамике.


Рассмотрим термодинамическую систему (ТДС), находящуюся в равновесном


состоянии. Предположим, что она совершает цикл Карно по пути


$T^1S^1T^2S^2T^1$, где $S^1 > S^2$, т.\,е. этот цикл однозначно 


характеризуется параметрами


$T^1$, $S^1$, $T^2$, $S^2$. Пусть $T^1$ и $T^2$ ---


температуры нагревателя и холодильника соответственно, причем $T^1 > T^2$.


Предположим, что в результате этого процесса ТДС от нагревателя получает


количество тепла $Q^1=T^1(S^1 - S^2)$ и отдает холодильнику количество


тепла $Q^2=T^2(S^2 - S^1)$. Допустим,


что $x$ --- состояние ТДС с энтропией $S^1=x^1$ и температурой $T^1 =


x^2$, а $y$ --- состояние системы с энтропией $S^2=y^1$ и температурой


$T^2=y^2$. Тогда интерпретируем двуметрическое пространство $M$ как пространство


состояний равновесной ТДС, а двуметрические функции $f^1(xy)$ и $f^2(xy)$


как количества тепла $Q^1$ и $|Q^2|$ соответственно ---


$$


f^1(xy)=Q^1=T^1(S^1 - S^2),\quad


f^2(xy)=|Q^2|=T^2(S^1 - S^2).


$$


Заметим, что пара


состояний $x$ и $y$ однозначно определяет цикл Карно


со следующими параметрами: $(T^1,S^1,T^2,S^2)$.


Согласно этой интерпретации инвариант (5) представляет собой коэффициент


полезного действия цикла Карно


$$


\eta=\frac{Q^1 - |Q^2|}{Q^1}=\frac{x^2 - y^2}{x^2}=\frac{T^1 - T^2}{T^1} < 1,


$$


а инвариант (6) --- полезная работа цикла Карно


$$


A=Q^1 - |Q^2|=(x^1 - y^1)(x^2 - y^2)=(S^1 - S^2)(T^1 - T^2).


$$


По своему физическому смыслу координаты $x^1$ и


$y^1$ строго положительны [2].





Рассмотрим движение $g$ плоскости $M$, которое произвольные


состояния $x$ и $y$ переводит в состояния $x'$ и $y'$ соответственно, т.\,е.


$S^{\prime 1}=g^1(S^1,T^1)$, $T^{\prime 1}=g^2(S^1,T^1)$, $S^{\prime 2} =


g^1(S^2,T^2)$, $T^{\prime 2}=g^2(S^2,T^2)$. Это приводит к тому, что


цикл Карно с параметрами $(S^1,T^1,S^2,T^2)$ переходит в цикл Карно с


параметрами $(S^{\prime 1},T^{\prime 1},S^{\prime 2},T^{\prime 2})$. Физически


это означает, что отображение $g$ устанавливает связь между двумя круговыми


процессами (циклами).


Следовательно, КПД $\eta$ и полезная работа $A$ этих циклов равны. Заметим,


что в этом случае


$$


S^{\prime}=e^{\alpha} S+b,\quad T^{\prime}=e^{-\alpha} T,


$$


т.\,е. температуры нагревателя и холодильника при преобразовании $g$ меняются


пропорционально с одним и тем же коэффициентом.


Выражая $e^{-\alpha}$ из преобразований температуры нагревателя и подставляя в


преобразования температуры холодильника этих циклов, получим


\begin{equation}


\frac{T^{\prime 1}}{T^1}=\frac{T^{\prime 2}}{T^2}.


\end{equation}


Для определения коэффициента $e^{-\alpha}$


рассмотрим состояние с эталонной температурой $T_0$, которое


при преобразовании $g$ переходит в состояние с температурой $T$.


Тогда $\alpha=\ln(T_0/T)$.


Поэтому согласно (7) температуры холодильников циклов Карно относительно


преобразований $g$ связаны равенством


$$


T^{\prime 2}=\frac{T}{T_0}T^2,


$$


где $T_0$ --- температура нагревателя первого из них.





Траекториями однопараметрической подгруппы группы (3)


\begin{equation}


S^{\prime}=e^{\alpha} S,\quad T^{\prime}=e^{-\alpha} T


\tag{$3'$}


\end{equation}


являются гиперболы, а также


вертикальные и горизонтальные прямые. Пусть группа


преобразований (3) переводит состояние $x$ в состояние $x'$. Тогда


на возрастание температуры система реагирует убыванием энтропии


и наоборот. Из (3) следует равенство


$S^{\prime}T^{\prime}=ST$.





Из инвариантности двуметрической функции $(f^1,f^2)$


пространства $M$ относительно группы движений следует, что цикл


Карно любой равновесной ТДС можно заменить на физически более выгодный


цикл Карно с тем же КПД $\eta$ и той же полезной работой $A$.


\medskip








{\bf 2. Метрические функции двуметрических двумерных многообразий.}


В данном параграфе определяются двуметрические многообразия и


строятся их метрические функции.





Пусть $M$ --- гладкое двумерное многообразие, $T(M)$ --- касательное


расслоение, каждый слой $T_x(M)$ которого будем рассматривать как


аффинное пространство.


Стандартный слой $R^2$ является линейным пространством, который также будем


рассматривать как аффинное пространство. 


Пусть $A(M)$ --- расслоение аффинных реперов $\wt{u}$ со структурной


группой $A(2,R)$, которая является группой аффинных преобразований.


Аффинный репер многообразия $M$ в точке $x$ состоит из


точки $p \in T_x(M)$ и линейного репера $u=(X_1,X_2)$ в $x$. Тогда


$\wt{u}=(p;X_1,X_2)$.


Заметим, что элементы из $A(2,R)$ можно представлять в виде


$\wt{a} =\left(\begin{smallmatrix}


a&\xi\\ 0&1\end{smallmatrix}\right)$,


причем $a \in GL(2,R)$, $\xi \in R^2$.


Каждый аффинный репер $\wt{u}$ из $A(M)$


можно рассмотреть как аффинное преобразование $R^2$ на $T_x(M)$ [3].


Пусть $o$ --- начало в $R^2$ и $(e_1,e_2)$ --- фиксированный базис в $R^2$.


Тогда $\wt{u}(o;e_1,e_2)=(p;X_1,X_2)$.





Рассмотрим гладкое отображение


\begin{equation}


\omega: A_x(M)\times A_x(M)\times T_x(M) \to R^2,


\end{equation}


которое задается следующей формулой. Возьмем произвольные


аффинные реперы $\wt{u},\wt{v} \in A_x(M)$.


Пусть произвольный вектор $X \in T_x(M)$ в


базисе $\wt{v}$ имеет координаты $X^1,\ X^2$.


Тогда для функции (8) имеем


\begin{equation}


\omega(\wt{u},\wt{v},X)=a_i^jX^ie_j+\xi^ie_i,\quad i,j=1,2.


\end{equation}


Предположим, что $\wt{v}=(q;Y_1,Y_2)$ и $\wt{u}=(p;X_1,X_2)$.


В отношении отображения (8) выдвинем требование $\wt{a}=\wt{X}^{-1}$,


где $\wt{X}$ --- матрица отличия репера $\wt{u}$ от репера $\wt{v}$.


Для открытого покрытия


$U_{\alpha}$ многообразия $M$ (в силу локальной


тривиальности расслоения $A(M)$) существует семейство изоморфизмов


$\varkappa_{\alpha}: A(U_{\alpha}) \to U_{\alpha}\times A(2,R)$,


определяющих сечения через $\wt{u}$ и $\wt{v}$ для


$\sigma_{\alpha}: U_{\alpha} \to A(M)$ по формулам


$\sigma_{1\alpha}(x)=\varkappa^{-1}_{1\alpha}(x,e)$ и $\sigma_{2\alpha}(x) =


\varkappa^{-1}_{2\alpha}(x,e)$ соответственно. Тогда соответствие


$x \mapsto \omega$


является гладким в произвольной координатной окрестности


$U \subset M$. Поэтому элементы матрицы $\wt{a}$ являются


гладкими функциями в $U$, которые будем называть


{\it структурными функциями}


двуметрического многообразия $M$.





\begin{lemnonum}


При переходе от системы координат $U$ к системе координат


$U'$ в точках их пересечения


структурные функции $\wt{a}$ преобразуются по закону


$$


a_{i}^{\prime j}=a_k^j\frac{\partial x^k}{\partial x^{'i}},\quad


\xi^{\prime i}=\xi^i,\quad i,j,k=1,2,


$$


где $a$ --- линейная компонента структурных функций $\wt{a}$ в


координатной окрестности $U$, а $a'$ --- линейная компонента матрицы


$\wt{a}{}'$ в координатной окрестности $U'$.


\end{lemnonum}





Из этой леммы следует, что функции (9) инвариантны относительно


произвольной замены координат.





Рассмотрим замкнутую подгруппу Ли $G \subset A(2,R)$.


Обозначим через $Q(M,G)$ или просто $Q$


редуцированное подрасслоение расслоения $A(M)=P(M,A(2,R))$.


Рассмотрим сужение отображения (8) на $Q_x$ относительно первого аргумента


\begin{equation}


\omega: Q_x(M)\times A_x(M)\times T_x(M) \to R^2.


\end{equation}


Это отображение представим в виде $\omega(\pi(\wt{u}),


\wt{v},X)$, где $\wt{u},\wt{v} \in A_x(M)$, а $\pi: A_x(M) \to


Q_x(M)$ --- гладкая проекция, которая реперы $b\psi(\wt{u}')$,


$b\in A(2,R)/G$, отображает в $\wt{u}{}'$.


Из гладкости отображения $\pi$ следует, что функция (10) является гладкой.


Тогда для фиксированного репера $\wt{v} \in A_x(M)$


\begin{equation}


\omega_{\Tilde{v}}(\wt{u},X)=a_i^jX^ie_j+\xi^je_j,\quad i,j=1,2,


\end{equation}


где $\wt{u}$ --- произвольный репер из $Q_x$. Поскольку


в (11) $\wt{u}$ --- произвольный репер


из $Q_x$, а $\wt{v}$ --- фиксированный репер из $A_x(M)$,


то для структурных функций $\wt{a}$ \ $G$-пространства справедливо разложение


\begin{equation}


\wt{a}=\wt{b}\wt{c},


\end{equation}


где в произвольной точке из $U$ матрица $\wt{b}$ --- произвольный элемент 


подгруппы $G$, а $\wt{c}$ --- некоторая


фиксированная матрица из $A(2,R)$. Ниже под $\wt{v}$


будет пониматься репер $(0;\partial_{x^1},\partial_{x^2})$.


Можно показать, что соответствие $x \mapsto \omega$ гладко.





Пусть подгруппа Ли $G$ как группа преобразований


задается уравнениями (4).


Тогда для структурных функций $\wt{a}$ имеем разложение (12), причем


в каждой точке из $U$


$$


\wt{b} =


\begin{pmatrix}


e^{\varepsilon\alpha}&0&b\\


0&e^{-\varepsilon\alpha}&(1-\varepsilon)a\\


0&0&1


\end{pmatrix}


,\quad \varepsilon=0,1.


$$





Поскольку $\wt{u}$ --- произвольный аффинный репер из $Q(M,G)$, то


отображение (11) можно представить в виде


\begin{equation}


\omega_{\Tilde{v}}(\wt{u},X)=((a_i^1X^i+\xi^1)e^{\varepsilon\alpha} +


b)e_1+((a_i^2X^i+\xi^2)e^{-\varepsilon\alpha}+(1-\varepsilon)a)e_2.


\tag{$11'$}


\end{equation}


Возьмем два вектора $X$ и $Y$ из пространства $T_x(M)$ и


определим с помощью отображения $(11')$ их образы в пространстве $R^2$:


$\omega_{\Tilde{v}}(\wt{u},X)$ и $\omega_{\Tilde{v}}(\wt{u},Y)$.


В силу произвольности выбора репера $\wt{u}$ получим два семейства


векторов в $R^2$, которые являются орбитами группы $G$.


Так как группа~$G$ двухпараметрическая и действует просто транзитивно в


$R^2$, то существуют два независимых двухточечных инварианта


этой группы --- (1) и (2). Поэтому 


приходим к двуметрическим функциям


\begin{gather}


f^1(X,Y)=a_i^1X^i - a_i^1Y^i,\quad


f^2(X,Y)=a_i^2X^i - a_i^2Y^i;\\


\begin{split}


f^1(X,Y)&=(a_i^1X^i - a_i^1Y^i)(a_i^2X^i+\xi^2),\\


f^2(X,Y)&=(a_i^1X^i - a_i^1Y^i)(a_i^2Y^i+\xi^2),


\end{split}


\end{gather}


где $i,j=1,2$.


Заметим, что двуметрические функции (13) и (14) определяются двумя


ковекторными полями и одним скалярным полем.





{\it Двумерное многообразие $M$ будем называть


двуметрическим, если в нем определена двуметрическая функция


$(13)$ или $(14)$.}


Двуметрическая функция инвариантна относительно


преобразований структурных функций $\wt{a} \to \wt{b}\wt{a}$.





Воспользуемся тем, что дифференциалы координат $dx$ преобразуются по векторному


закону, т.\,е. образуют вектор. Тогда, полагая в (13) и (14) $Y^i =


0$, $X^i=dx^i$, получим


\begin{gather}


f^1=a_i^1dx^i,\quad f^2=a_i^2dx^i; \tag{$13'$} \\


f^1=(a_i^1dx^i)(a_i^2dx^i+\xi^2),\quad


f^2=a_i^1dx^i\xi^2.\tag{$14'$}


\end{gather}


Если для некоторых двуметрических пространств в $U$


можно подобрать такую матрицу $\wt{b}$, что


$a_i^j=\delta_i^j$, $\xi^i=x^i$, то метрические


функции приобретут вид $(1')$ и $(2')$.


Такие пространства называются локально плоскими.


\medskip











{\bf 3. Связности двуметрических двумерных пространств.}


Теперь приступим к исследованию связностей в двуметрических 


двумерных пространствах. Предположим, что в расслоении аффинных 


реперов $L(M)$ задана аффинная связность.





Аффинная связность в $L(M)$ двуметрического двумерного пространства $M$ 


называется {\it согласованной}, если параллельный перенос слоев


из $T(M)$ сохраняет двуметрическую функцию $f=(f^1,f^2)$.


Из определения согласованной связности пространства $M$


следуют равенства


\begin{equation}


\nabla_kf^1(X,Y)=0,\quad \nabla_kf^2(X,Y)=0.


\end{equation}








\begin{thm}


Компоненты символов Кристоффеля $\Gamma$ согласованной связности в


координатной окрестности $U$ двуметрических


двумерных пространств с метрическими функциями $(13)$ и $(14)$ имеют следующие


выражения\/${:}$





для двуметрики $(13)$


\begin{equation}


\Gamma^l_{ik}=\ov{a}{}_j^l\frac{\partial a_i^j}{\partial x^k},


\end{equation}


где $\ov{a}$ --- матрица, обратная к матрице $a;$





для двуметрики $(14)$


\begin{equation}


\Gamma^n_{ik}= \ov{a}{}^n_m\frac{\partial a^m_i}{\partial x^k}+


\ov{a}{}^n_m\wt{a}{}^m_i\frac{\partial ln\xi^2}{\partial x^k},\quad


i,j,k,l,n,m=1,2,


\end{equation}


причем $\wt{a}{}^m_i=\varepsilon(m)a^m_i$,


$\varepsilon(1)=1$, $\varepsilon(2)=-1$.


\end{thm}





\begin{pf}


Предположим, что


$X$ --- параллельно переносимое векторное поле. Вычисляя ковариантные


производные метрических функций $f^1$ и $f^2$ и приравнивая согласно (15)


их к нулю, для первой двуметрики получаем


$$


\nabla_ka^1_i=0,\quad \nabla_ka^2_i=0.


$$


Расписывая эти две ковариантные производные, а затем их объединяя и выражая


символы $\Gamma$, приходим к выражениям (16).


Во втором случае вычисляя ковариантные производные, имеем


\begin{gather*}


\bigg(\frac{\partial a^1_i}{\partial x^k} - \Gamma^{l}_{ik}


a^1_l\bigg)(a^2_sX^s+\xi^2)+a^1_i\bigg(\bigg(\frac{\partial a^2_s}{\partial


x^k}-\Gamma^{l}_{sk}a^2_l\bigg)X^s+\frac{\partial \xi^2}{\partial


x^k}\bigg)=0, \\


%


\bigg(\frac{\partial a^1_i}{\partial x^k} - \Gamma^{l}_{ik}


a^1_l\bigg)(a^2_sY^s+\xi^2)+a^1_i\bigg(\bigg(\frac{\partial a^2_s}{\partial


x^k} - \Gamma^{l}_{sk}a^2_l\bigg)Y^s+\frac{\partial \xi^2}{\partial


x^k}\bigg)=0.


\end{gather*}


В касательном пространстве $T_x(M)$, где $x$ --- произвольная точка из


$M$, эти выражения представляют собой линейные


функции относительно координат векторов $X$ и $Y$.


Из этих равенств следует, что их


коэффициенты равны нулю. Поэтому приходим к (17).


\end{pf}








\begin{cornonum}


Символы Кристоффеля (16) инвариантны относительно


тождественных преобразований структурных функций $a$, а символы


Кристоффеля (17) инвариантны относительно преобразований


$a^1_i\to e^{\alpha}a^1_i$,


$a^2_i\to e^{-\alpha}a^2_i$, $\xi^2\to e^{-\alpha}\xi^2$,


$\xi^1\to e^{\alpha}\xi^1+b$.


\end{cornonum}





Найдем значения компонент символов Кристоффеля согласованной связности


локально плоских пространств.


Из выражений (16) будет следовать


\begin{equation}


\Gamma^k_{ij}=0,\tag{$16'$}


\end{equation}


а из выражений $(17)$ ---


\begin{equation}


\Gamma^1_{1k}=\frac{\partial (ln\xi^2)}{\partial x^k},\ \ 


\Gamma^2_{2k}=-\frac{\partial (ln\xi^2)}{\partial x^k},\ \ 


\Gamma^i_{jk}=0,\ \ i\neq j.\tag{$17'$}


\end{equation}





Для нахождения компонент тензора кривизны $R$ (в выше рассмотренной координатной


окрестности $U$) связности с символами


Кристоффеля $(16')$ и $(17')$ воспользуемся формулой


$$


R_{jkl}^i=\frac{\partial \Gamma _{lj}^i}{\partial x^k}+\Gamma _{lj}^m\Gamma


_{km}^i-\frac{\partial \Gamma _{kj}^i}{\partial x^l}-\Gamma _{kj}^m\Gamma


_{lm}^i.


$$


Прямой подстановкой несложно убедиться в том, что компоненты тензора кривизны


$R$ для этих связностей тождественно обращаются в нуль.





Вычислим компоненты тензора кручения $T$,


$T^k_{ij}=\Gamma^k_{ij} - \Gamma^k_{ji}$.


Воспользовавшись выражениями (16) и (17), находим





для многообразия с двуметрикой (13)


$$


T^k_{ij}=\ov{a}{}_l^k\bigg(\frac{\partial a_i^l}{\partial x^j} -


\frac{\partial a_j^l}{\partial x^i}\bigg);


$$





для многообразия с двуметрикой (14)


$$


T^k_{ij}=\ov{a}{}_l^k\bigg(\frac{\partial a_i^l}{\partial x^j} -


\frac{\partial a_j^l}{\partial x^i}\bigg) +


\ov{a}{}_l^k\bigg(\wt{a}{}^l_i\frac{\partial \ln \xi^2}{\partial x^j} -


\wt{a}{}^l_j\frac{\partial \ln \xi^2}{\partial x^i}\bigg),


$$


причем $i,j,k,l=1,2$.





Кривая $x_t$, $a<t<b$, двуметрического двумерного многообразия $M$ с


согласованной


связностью называется {\it геодезической}, если касательное векторное поле $X


=\Dot{x}_t$, определенное вдоль $x_t$, параллельно вдоль $x_t$, т.\,е.


$\nabla_XX$ существует и равно нулю для всех $t$ [3].





Из определения геодезической следует, что она удовлетворяет системе


уравнений


\begin{equation}


\ddot{x}^l+\Gamma _{ij}^l\Dot{x}^i\Dot{x}^j=0,\quad i,j,l=1,2,


\end{equation}


где точка над функцией означает производную по параметру $t$.


Наша ближайшая цель состоит в том, чтобы уравнения геодезических (18)


для двуметрических двумерных пространств проинтегрировать полностью один раз


в нормальной координатной окрестности.





\begin{thm} 


Геодезическая линия для пространства с двуметрикой $(13)$


в координатной окрестности $U$ относительно согласованной


связности является решением системы дифференциальных уравнений первого порядка


\begin{equation}


a^1_i\dot{x}^i=a^1,\quad a^2_j\dot{x}^j=a^2,\quad i,j=1,2,


\end{equation}


где $a^1,a^2=\const$.


\end{thm}





\begin{pf}


Подставляя выражения для


символов Кристоффеля (16) в уравнения (18), а затем


умножая слева на матрицу $a^n_l$ и суммируя, приходим


к уравнениям


$a^n_l\ddot{x}^l+ \Dot{a}_l^n\dot{x}^l=0$.


Интегрируя эти уравнения, получаем (19).


\end{pf}





\begin{thm}


Геодезическая линия для пространства с двуметрикой $(14)$


в координатной окрестности $U$ относительно согласованной


связности является решением системы дифференциальных уравнений первого порядка


\begin{equation}


a^1_i\dot{x}^i=a^1/\xi^2,\quad a^2_j\dot{x}^j=a^2\xi^2,\quad i,j=1,2,


\end{equation}


где $a^1,a^2=\const$, причем функция $\xi^2$ в $U$ отлична от нуля.


\end{thm}





\begin{pf}


Подставляя выражения


для символов Кристоффеля (17) в уравнения (18), а затем


умножая слева на матрицу $a^n_l$ и суммируя, приходим


к уравнениям


$$


a^1_l\ddot{x}^l+ \Dot{a}_l^1\Dot{x}^l +


a^1_l\Dot{x}^l\Dot{\xi}^2/\xi^2=0,\quad


a^2_l\ddot{x}^l+ \Dot{a}_l^2\Dot{x}^l -


a^1_2\Dot{x}^l\Dot{\xi}^2/\xi^2=0.


$$


Интегрируя эти уравнения, получаем (20).


\end{pf}





Пусть $K_U$ --- множество решений систем уравнений вида (19) в координатной


окрестности $U$, а $N_U$ --- множество решений систем уравнений вида (20).


Заметим, что в $K_U$ есть геодезическая линия пространства с двуметрикой (13),


а в $N_U$ --- геодезическая линия пространства с двуметрической функцией (14).


Обратим внимание на то, что произвольная кривая из семейства $K_U$


является интегральной кривой векторного поля $\lambda$ в $U$, для которого


\begin{equation}


\lambda^1=\ov{a}^1_ja^j,\quad \lambda^2=\ov{a}^2_ja^j,


\end{equation}


где $\ov{a}$ --- матрица, обратная к матрице $a$.


Эти интегральные кривые являются решениями системы уравнений


\begin{equation}


\Dot{x}^1=\ov{a}{}^1_ja^j,\quad 


\Dot{x}^2=\ov{a}{}^2_ja^j.\tag{$19'$}


\end{equation}


Аналогично, кривая из семейства $N_U$


является интегральной кривой векторного поля в $U$


$$


\lambda^1=\ov{a}{}^1_1a^1/\xi^2+\ov{a}{}^1_2a^2\xi^2,\quad


\lambda^2=\ov{a}{}^2_1a^1/\xi^2+\ov{a}{}^2_2a^2\xi^2,


$$


которая является решением системы уравнений


\begin{equation}


\Dot{x}^1=\ov{a}{}^1_1a^1/\xi^2+\ov{a}{}^1_2a^2\xi^2,\quad


\Dot{x}^2=\ov{a}{}^2_1a^1/\xi^2+\ov{a}{}^2_2a^2\xi^2.\tag{$20'$}


\end{equation}





С данного момента и до конца параграфа под $U$ будет пониматься


нормальная координатная окрестность с началом в точке $x_0$.


Произвольную точку нормальной


окрестности можно соединить с началом единственной геодезической.


Обозначим через $K_U(x_0) \subset K_U$ ($N_U(x_0) \subset N_U$)


подмножество кривых, проходящих через точку $x_0$.


Для выяснения того, какая кривая в $K_U(x_0)$ или в $N_U(x_0)$


является геодезической,


воспользуемся определением, т.\,е. тем,


что при параллельном переносе


касательный вектор вдоль геодезической переходит


также в касательный вектор. Тогда


\begin{equation}


\frac{\partial \lambda^i}{\partial x^k}\lambda^k +


\Gamma^i_{jk}\lambda^j\lambda^k=0,\quad i,j,k=1,2.


\end{equation}


Эта система уравнений в $U$ всегда интегрируема,


поскольку через точку $x_0$ в произвольном направлении выходит


единственная геодезическая линия. Расписывая (22), получаем


\begin{gather*}


\frac{\partial \lambda^1}{\partial x^1}\lambda^1 +


\frac{\partial \lambda^1}{\partial x^2}\lambda^2 +


\Gamma^1_{11}\lambda^1\lambda^1 +


\Gamma^1_{12}\lambda^1\lambda^2 +


\Gamma^1_{21}\lambda^2\lambda^1 +


\Gamma^1_{22}\lambda^2\lambda^2=0, \\


\frac{\partial \lambda^2}{\partial x^1}\lambda^1 +


\frac{\partial \lambda^2}{\partial x^2}\lambda^2 +


\Gamma^2_{11}\lambda^1\lambda^1 +


\Gamma^2_{12}\lambda^1\lambda^2 +


\Gamma^2_{21}\lambda^2\lambda^1 +


\Gamma^2_{22}\lambda^2\lambda^2=0.


\end{gather*}





Для многообразия с двуметрикой (13)


подставляя в (22) соответствующие выражения для векторного


поля (21) и выражения для символов Кристоффеля (16),


приходим к тождеству.


Таким образом, в пространстве с двуметрической функцией (13) в


нормальной координатной окрестности $U$ с началом в точке $x_0$


геодезической, проходящей через $x_0$, является произвольная кривая из семейства


$K_U(x_0)$.


Аналогичным образом в пространстве с двуметрикой (14)


для геодезической из $N_U(x_0)$


выполняется тождество


\begin{equation}


\frac{\partial b^i}{\partial x^k}\ov{a}{}^k_sb^s +


\wt{a}{}^i_j\frac{\partial


\ln\xi^2}{\partial x^k}\ov{a}{}^j_mb^m\ov{a}{}^k_sb^s=0,


\end{equation}


причем $b^i=a^i/(\xi^2)^{\varepsilon(i)}$, $i,j,k,s,m=1,2$.


Структурные функции в окрестности $U$ с началом в $x_0$,


произвольные для пространства с двуметрикой (13) и


удовлетворяющие условию (23) для пространства с двуметрикой


(14), будем называть {\it каноническими}.


Если для пространства


с двуметрикой (14) в нормальной окрестности $U$ справедливы равенства


$a^i_j=\delta^i_j$, то условия (23)


выполняются тождественно, поэтому каноническими структурными


функциями являются $\delta^i_j$, $\xi^1$, $\xi^2$.





Найдем геодезические линии некоторых пространств.


Возьмем пространство с двуметрикой (13), которое в нормальной окрестности


$U$ с началом в $x_0$ является плоским, т.\,е. в $U$ \ $a^i_j=\delta^i_j$.


Тогда в этой окрестности для уравнений геодезических линий


$(19')$ с начальными условиями


$x^1(0)=x^1_0=0$, $x^2(0)=x^2_0=0$ имеем


$$


x^1(t)=a^1t,\quad x^2(t)=a^2t,


$$


т.\,е. геодезическими, проходящими через $x_0$, являются прямые,


выходящие из этой точки.





Аналогично, для пространства с двуметрикой (14), для которого


в нормальной окрестности $U$ с началом в точке $x_0$ \ $a^i_j=\delta^i_j$,


уравнения геодезических линий $(20')$, проходящих через точку $x^1(0)=x^1_0=0$,


$x^2(0)=x^2_0=1$, принимают вид


$$


\Dot{x}^1=a^1/\xi^2,\quad \Dot{x}^2=a^2\xi^2,


$$


причем условия (23) выполняются тождественно, т.\,е. любая кривая,


удовлетворяющая последней системе и проходящая через начало,


является геодезической.


Заметим, что векторные поля


$$


\lambda^1=a^1/\xi^2,\quad


\lambda^2=a^2\xi^2


$$


не имеют особенностей.





Рассмотрим случай локально плоского пространства,


когда в $U$ еще $\xi^i=x^i$. Тогда, интегрируя уравнения


геодезических с начальными условиями $x^1(0)=x^1_0=0$,


$x^2(0)=x^2_0=1$, получим


\begin{alignat*}{2}


x^1(t)&=a^1t, &\quad x^2(t)&=1, \\


x^1(t)&=- \frac{a^1}{a^2}e^{-a^2t}, &\quad 


x^2(t)&=e^{a^2t},


\end{alignat*}


причем $c_2\neq 0$, т.\,е. геодезическими являются прямые и гиперболы,


проходящие через точку $(0,1)$.


\medskip











{\bf 4. Двуметрические двумерные пространства, допускающие группы изометрий.}


Исследуем теперь двуметрические пространства $M$, допускающие группу изометрий.


Преобразование двуметрического двумерного многообразия $M$ называется 


изометрией, если


оно сохраняет метрическую функцию $f$. Пусть $X$ --- векторное поле в


двуметрическом многоообразии $M$, которое в координатной


окрестности $U$ произвольной точки $x$ определяет однопараметрическую группу


изометрий $\varphi_t$. Это векторное поле называется инфинитезимальной


изометрией двуметрического пространства $M$. При этом выполняются равенства


$$


L_Xf^1=0,\quad L_Xf^2=0,


$$


где $L_X$ --- производная Ли в направлении векторного поля $X$. Поэтому 


векторное поле $X$ является


инфинитезимальной изометрией.





\begin{thm}


Для того чтобы в координатной окрестности $U$ произвольной


точки двуметрического двумерного многообразия $M$ векторное 


поле $X$ было инфинитезимальной


изометрией, необходимо и достаточно, чтобы имела место следующая система


дифференциальных уравнений\/${:}$





для пространств с двуметрикой $(13')$


\begin{equation}


\frac{\partial a_j^i}{\partial x^k}X^k+a_k^i


\frac{\partial X^k}{\partial x^j}=0;


\end{equation}





для пространств с двуметрикой $(14')$


\begin{equation}


\frac{\partial a_j^i}{\partial x^k}X^k+a_k^i\frac{\partial X^k}{\partial x^j}


=\varepsilon(i)a^i_j\frac{\partial \ln\xi^2}{\partial x^k}X^k,


\end{equation}


где $i,j,k=1,2$, причем $X^1$ и $X^2$ --- координаты векторного поля $X$,


$\varepsilon(1)= -1$, $\varepsilon(2)=1$.


\end{thm}





\begin{pf}


Воспользуемся выражениями для


производной Ли $L_X$ в направлении векторного поля $X$ \ 1-форм


и функций.


Вычисляя производные Ли двуметрических функций $(13')$ и $(14')$, приходим к


следующим равенствам:





для двуметрики $(13')$


$$


L_Xf^1=\bigg( \frac{\partial a_j^1}{\partial x^k}X^k+a_k^1


\frac{\partial X^k}{\partial x^j}\bigg) dx^j,\quad


L_Xf^2=\bigg( \frac{\partial a_j^2}{\partial x^k}X^k+a_k^2


\frac{\partial X^k}{\partial x^j}\bigg) dx^j;


$$





для двуметрики $(14')$


\begin{align*}


L_Xf^1&=\bigg(\frac{\partial a_i^1}{\partial x^k}X^k+a_k^1


\frac{\partial X^k}{\partial x^i}\bigg) dx^i(a^2_jdx^j+\xi^2)+


a^1_idx^i\bigg[\bigg(


\frac{\partial a_j^2}{\partial x^k}X^k+ a_k^2\frac{\partial X^k


}{\partial x^j}\bigg) dx^j+\frac{\partial \xi^2}{\partial x^k}X^k\bigg],\\


%


L_Xf^2&=\bigg(\frac{\partial a_i^1}{\partial x^k}X^k+a_k^1\frac{\partial X^k


}{\partial x^i}\bigg) dx^i\xi^2+a^1_idx^i\frac{\partial \xi^2}{\partial


x^k}X^k.


\end{align*}


Из этих равенств следует, что если векторное поле $X$ является


инфинитизимальной изометрией, то имеют место уравнения (24) и (25)


соответственно, и наоборот.


\end{pf}





\begin{cornonum}


Уравнения (24) инвариантны относительно


тождественных преобразований структурных функций $a$, а уравнения


(25) инвариантны относительно преобразований $a^1_i\to e^{\alpha}a^1_i$,


$a^2_i\to e^{-\alpha}a^2_i$, $\xi^2\to e^{-\alpha}\xi^2$,


$\xi^1\to e^{\alpha}\xi^1+b$.


\end{cornonum}





Воспользовавшись, далее, выражениями для символов Кристоффеля (16) и (17),


после несложных преобразований из (24) и (25) приходим





для пространств с двуметрикой $(13')$


$$


\frac{\partial X^i}{\partial x^j}+\Gamma^i_{jk}X^k=0;


$$





для пространств с двуметрикой $(14')$


$$


\frac{\partial X^i}{\partial x^j}+\Gamma^i_{jk}X^k - 2G^i_{jk}X^k= 0,


$$


где введено сокращающее обозначение


$G^i_{jk}=\ov{a}{}^i_l\wt{a}{}^l_j\frac{\partial \ln\xi^2}{\partial


x^k}$.





Проинтегрируем уравнения (24) и (25) для частного случая, когда в


координатной окрестности $U$ произвольной точки $a^i_j=\delta^i_j$.





\begin{thm}


Операторы группы инфинитезимальных изометрий в


координатной окрестности $U$ с двуметрикой $(1')$ имеют вид


\begin{equation}


X^1=\partial_{x^1},\quad


X^2=\partial_{x^2}.


\end{equation}


\end{thm}





\begin{pf}


Задача сводится к интегрированию системы дифференциальных


уравнений


$$


\frac{\partial X^1}{\partial x^1}=0,\quad


\frac{\partial X^1}{\partial x^2}=0,\quad


\frac{\partial X^2}{\partial x^1}=0,\quad


\frac{\partial X^2}{\partial x^2}=0,


$$


от которой приходим к оператору алгебры Ли


$X=a^1\partial_{x^1}+a^2\partial_{x^2}$,


где $a^1$ и $a^2$ --- произвольные постоянные. Полагая в первом случае


$a^1=1$, $a^2=0$, а во втором $a^1=0$, $a^2=1$, приходим к


операторам (26).


\end{pf}





Несложно убедиться в том, что геодезическая линия пространства с двуметрикой


$(1')$ является орбитой группы изометрий с операторами (26).





Записывая уравнения (25) для случая $a^i_j=\delta^i_j$, получаем


$$


X^1=ax^1+c,\quad


X^2=-ax^2 - d,\quad


\frac{\partial X^1}{\partial x^1}\xi ^2+\frac{\partial \xi ^2}{\partial x^k}%


X^k=0.


$$


Далее, наша цель состоит в нахождении функций $\xi^2$, для которых последняя


система разрешима. Заметим, что функция $\xi^2 \neq \const$, поскольку иначе


от $(14')$ сможем перейти к $(13')$, для этого достаточно положить


$g^1=f^2/\xi^2$, $g^2=\xi^2f^1/f^2 - \xi^2$. Поэтому возможны два случая:


1)~$a\neq 0$, 2)~$a=0$, $c^2+d^2\neq0$.





В первом случае имеем уравнение для функции $\xi^2$


\begin{equation}


\xi^2+ \frac{\partial \xi ^2}{\partial x^1}(x^1+\alpha) -


\frac{\partial \xi^2}{\partial x^2}(x^2+\beta)=0,


\end{equation}


где $a,\ c,\ d$ --- произвольные постоянные, причем $\alpha=c/a$, $\beta=d/a$.





Во втором случае приходим к уравнению


\begin{equation}


c\frac{\partial \xi ^2}{\partial x^1} -


d\frac{\partial \xi^2}{\partial x^2}=0.


\end{equation}





Интегрируя уравнение (27), находим


\begin{equation}


\xi^2=\frac{\psi((x^1+\alpha)(x^2+\beta))}{x^1+\alpha}.


\end{equation}


Заметим, что функция (29) не зависит от переменной $x^1$ тогда и только


тогда, когда $\psi$ --- однородная функция, т.\,е.


\begin{equation}


\psi=\theta (x^1+\alpha)(x^2+\beta),\tag{$29'$}


\end{equation}


где $\theta=\const$.


Интегрируя уравнение (28) при $c\neq 0$, получим


\begin{equation}


\xi^2=\varphi(\gamma_1x^1+x^2),


\end{equation}


где введено обозначение $\gamma_1=d/c$.


Аналогично, если $d\neq 0$, то интегралом уравнения (28) будет


\begin{equation}


\xi^2=\varphi(x^1+\gamma_2x^2),


\end{equation}


где $\gamma_2=c/d$.





\begin{thm}


Двуметрическое пространство с двуметрикой


$$


f^1=dx^1(dx^2+\xi^2),\quad f^2=dx^1\xi^2


$$


в координатной окрестности $U$


допускает двупараметрическую группу изометрий, если


$\xi^2=\theta(x^2+\beta)$,


причем ее операторы $X_1=\partial_{x^1}$,


$X_2=x^1\partial_{x^1} - (x^2+\beta)\partial_{x^2};$


допускает однопараметрическую группу изометрий, если функция $\xi^2$


имеет вид либо $(29)$, либо $(30)$, либо $(31)$, причем $\psi$ ---


неоднородная функция, с операторами


$$


X_1=(x^1+\alpha)\partial_{x^1} - (x^2+\beta)\partial_{x^2};\quad


X_1=\partial_{x^1} - \gamma_1\partial_{x^2};\quad


X_1=\gamma_2\partial_{x^1} - \partial_{x^2},


$$


а в остальных случаях групп изометрий нет.


\end{thm}





\begin{cornonum}


Если в утверждении теоремы 6 предполагать, что $\xi^i =


x^i$ (локально плоский случай), то приходим к группе изометрий с базисными


векторными полями


\begin{equation}


X_1=\partial_{x^1},\quad


X_2=x^1\partial_{x^1} - x^2\partial_{x^2}.


\end{equation}


\end{cornonum}





Несложно установить, что геодезическая линия пространства с двуметрикой


$(2')$ в $U$ является орбитой группы изометрий с операторами (32).


\medskip








{\bf 5. Двуметрические пространства в неравновесной термодинамике.}


Двумерная геометрия с двуметрикой $(14')$ допускает содержательную интерпретацию


в неравновесной термодинамике.


Рассмотрим неравновесную ТДС и выделим в ней такую


область, в которой можно говорить о температуре и об энтропии.


Пусть $M$ --- пространство состояний этой ТДС.


Предположим, что система совершает цикл Карно по пути $T^1S^1T^2S^2T^1$,


причем предполагается, что $T^1$, $T^2$ и $S^1$, $S^2$ мало отличаются


друг от друга соответственно.


Итак, цикл Карно однозначно характеризуется параметрами


($T^1$, $S^1$, $T^2$, $S^2$). Пусть $T^1$ и $T^2$ --- температуры нагревателя


и холодильника соответственно, причем $T^1 > T^2$, $S^1>S^2$.


Предположим, что в результате этого процесса система от нагревателя получает


количество тепла $Q^1$ и отдает холодильнику количество тепла $Q^2$.


Интерпретируем двуметрическое пространство $M$ как пространство


состояний неравновесной ТДС, а двуметрические функции $f^1$ и $f^2$


как количества тепла $Q^1$ и $|Q^2|$ соответственно.


Пусть $U$ --- координатная окрестность в $M$, т.\,е. множество близких


состояний. Рассмотрим два состояния из $U$: $1=(x^1+dx^1,x^2+dx^2)$ и


$2=(x^1,x^2)$. Тогда интерпретируем температуру и энтропию


этих состояний так: $T^1=x^2+dx^2$, $S^1=x^1+dx^1$ и


$T^2=x^2$, $S^2=x^1$. Заметим, что пара


состояний $1$ и $2$ однозначно определяет цикл Карно с параметрами


$(T^1,S^1,T^2,S^2)$. Введем функции


\begin{equation}


A^1_T=\xi^2+a^2_idx^i,\ \ A_S^1=\xi^1+a^1_idx^i,


\ \ A_T^2=\xi^2,\ \ A_S^2=\xi^1.


\end{equation}


Структурные функции


$a^i_j=\delta^i_j$ и $\xi^i=x^i$ в $U$ соответствуют


локально равновесному состоянию ТДС. Из (33) следует,


что в локально равновесном состоянии $A^1_T=T^1$, $A^1_S=S^1$,


$A^2_T=T^2$, $A^2_S=S^2$. Обозначим


$$


\delta A_S=A_S^1 - A_S^2=a^1_idx^i,\ \ 


\delta A_T=A_T^1 - A_T^2=a^2_idx^i.


$$


Эти выражения в локально плоском случае являются


полными дифференциалами и совпадают соответственно с $dS$ и $dT$.


Поскольку $M$ является двумерным многообразием, то


эти формы обладают интегрирующими множителями $b_S$ и $b_T$ соответственно.


Наличие этих интегрирующих множителей свидетельствует о том, что


в неравновесном состоянии в ТДС


функционируют ``источники энтропии'', которые при переходе в


равновесное состояние иссякают. Тогда


$$


\delta A_S=\frac{1}{b_S}dB_S,\quad 


\delta A_T=\frac{1}{b_T}dB_T,


$$


где $B_S$ и $B_T$ --- соответствующие интегралы, причем в равновесном


состоянии $b_S=1$, $b_T=1$. Таким образом,


$$


f^1=(\delta A_T+A_T)\delta A_S,\quad


f^2=A_T\delta A_S,


$$


где $A_T=A_T^2$. Итак, количество тепла для неравновесной ТДС


$\delta Q=A_T\delta A_S$.


Заметим, что в равновесных состояниях эта формула совпадает с


$\delta Q=TdS$. Поэтому согласно второму началу термодинамики и


нашей интерпретации


\begin{equation}


\delta Q=A_T\delta A_S=\frac{A_T}{b_S}dB_S=\xi^2a^1_idx^i \le x^2dx^1.


\end{equation}


Заметим, что в случае равновесных состояний (34) представляет собой равенство.


Таким образом, для неравновесных процессов первое начало термодинамики


записывается в виде


$$


dU=\frac{A_T}{b_S}dB_S - PdV.


$$





Согласно нашей интерпретации функция (5) представляет собой коэффициент


полезного действия цикла Карно


\begin{equation}


\eta=\frac{a^2_idx^i}{\xi^2+a^2_jdx^j}=\frac{\delta A_T}{A_T +


\delta A_T} < 1, \tag{$5'$}


\end{equation}


а функция (6) --- полезная работа цикла Карно


$$


A=a^1_idx^ia^2_jdx^j=\delta A_S\delta A_T =


\frac{dB_S}{b_S}\frac{dB_T}{b_T}.


$$


Из неравенства $(5')$ следует, что $\xi^2 > 0$, $A_T > 0$.


По теореме Карно [2] КПД любого теплового цикла с температурой


нагревателя $T^1$ и температурой холодильника $T^2$ меньше либо равна


КПД равновесного цикла Карно с теми же $T^1$ и $T^2$. Поэтому


$$


\frac{a^2_idx^i}{\xi^2+a^2_jdx^j} \le


\frac{T^1 - T^2}{T^1}=\frac{dx^2}{x^2+dx^2}.


$$


Заметим, что в равновесном состоянии последнее неравенство


превращается в равенство.











\begin{thebibliography}{9}





\bibitem{1}


Михайличенко Г.Г. {\em Полиметрические геометрии}. -- Новосибирск, 2001. 


-- 144 c.


\bibitem{2}


Румер Ю.Б., Рывкин М.Ш. {\em Термодинамика, статистическая физика и кинетика}.


-- М.: Наука, 1977. -- 552 c.


\bibitem{3}


Кобаяси Ш., Номидзу К. {\em Основы дифференциальной геометрии}. 


-- М.: Наука, 1981. -- 344 c.





\end{thebibliography}





\vskip 2cm





\post{\it Горно-Алтайский государственный}{\it Поступила}


\post{\it университет}{06.05.2003}





\label{end}


\end{document}


